% from ega4-20-fr.tex
\oldpage[IV-4]{241}
$M_\alpha\subset U_\alpha$, donc $M_\lambda\subset U_\lambda$ pour $\lambda\geq\alpha$ et $M\subset U$ en vertu de \sref{IV.20.3.8.3}, ce qui prouve la première assertion.
Inversement, supposons $M$ fini et $U$ dense dans $S$, donc $M\subset U$ ; notons que les $U_\lambda$ sont ouverts, donc ind-constructibles et les $M_\lambda$ finis, donc pro-constructibles \sref{IV.1.9.6}.
La seconde assertion résulte alors de \sref{IV.8.3.2}.

\phantomsection\label{IV.20.3.8.5}
\emph{20.3.8.5. Fin de la démonstration de \sref{IV.20.3.8.2}.}
Supposons la condition~(ii) vérifiée, et notons que $X$ est alors \emph{réduit} en vertu de \sref{IV.8.7.1} ; il revient donc au même de dire que $U_\lambda$ (resp. $U$) est schématiquement dense dans $X_\lambda$ (resp. $X$) ou qu'il est dense dans $X_\lambda$ (resp. $X$), et la conclusion résulte de \sref{IV.20.3.8.4} appliqué aux $X_\lambda$ et à $X$.

\phantomsection\label{IV.20.3.8.6}
\emph{20.3.8.6. Fin de la démonstration de \sref{IV.20.3.8}.}
Pour montrer que l'application \sref{IV.20.3.8.1} est injective, considérons deux morphismes $u'_\lambda$, $u''_\lambda$ d'un ouvert schématiquement dense $U_\lambda\subset X_\lambda$ dans $Y_\lambda$, et supposons que leurs images $u'$, $u''$, morphismes de $U$ dans $Y$, coïncident dans un ouvert schématiquement dense $V$ de $U$.
En outre, dans l'une ou l'autre des hypothèses~(i),~(ii), on peut supposer $V$ \emph{quasi-compact} ; c'est évident si $X$ est noethérien ; sinon, comme $X$ n'a qu'un nombre fini de points maximaux et est réduit, il suffit de remplacer $V$ par la réunion d'un nombre fini de voisinages ouverts affines de ces points maximaux (contenus dans $V$ par hypothèse).
Alors il existe un $\lambda\geq\alpha$ tel que $V$ soit l'image réciproque d'un ouvert quasi-compact $V_\lambda$ de $X_\lambda$ \sref{IV.8.2.11}, et il résulte de \sref{IV.20.3.8.2} qu'en prenant $\lambda$ assez grand, on peut supposer que $V_\lambda$ est schématiquement dense dans $X_\lambda$.
Il résulte alors de \sref{IV.8.8.2}[(i)] qu'en prenant $\lambda$ assez grand, $u'_\lambda$ et $u''_\lambda$ coïncident dans $V_\lambda$, donc appartiennent au même pseudo-homomorphisme.

Prouvons enfin que l'application \sref{IV.20.3.8.1} est surjective.
On considère maintenant un morphisme $u$ d'un ouvert schématiquement dense $U$ de $X$, dans $Y$ ; comme ci-dessus, on peut supposer $U$ quasi-compact et quasi-séparé ($U$ pouvant être remplacé, dans le cas~(ii), par une réunion d'un nombre fini d'ouverts affines deux à deux disjoints).
On peut donc encore supposer qu'il existe un $\lambda\geq\alpha$ tel que $U$ soit l'image réciproque d'un ouvert quasi-compact $U_\lambda$ de $X_\lambda$ \sref{IV.8.2.11} qui est automatiquement quasi-séparé, et par \sref{IV.20.3.8.2} on peut supposer en outre que $U_\lambda$ soit schématiquement dense dans $X_\lambda$.
Puisque les $Y_\lambda$ sont supposés localement de présentation finie, il résulte de \sref{IV.8.8.2}[(i)] qu'il existe un $\lambda$ tel que $u$ soit l'image d'un morphisme $u_\lambda:U_\lambda\to Y_\lambda$ ; d'où la conclusion.

\begin{remarks}[20.3.8.7]
\label{IV.20.3.8.7}
\begin{enumerate}[label=(\roman*)]
  \item Pour prouver que l'application \sref{IV.20.3.8.1} est \emph{injective}, il n'est pas nécessaire, dans l'hypothèse~(i) de \sref{IV.20.3.8}, de supposer $X$ noethérien.
  En effet, le lemme \sref{IV.20.3.8.2} n'utilise pas cette hypothèse.
  Avec les notations de \sref{IV.20.3.8.6}, soit $Z_\lambda$ le sous-préschéma des coïncidences de $u'_\lambda$ et $u''_\lambda$, et soit $Z$ le sous-préschéma des coïncidences de $u'$ et $u''$ ; il résulte de la définition \sref{IV.17.4.5} et de \sref[I]{I.3.3.10.1} que $Z$ est limite projective des $Z_\mu$ pour $\mu\geq\lambda$.
  Or, par hypothèse, $Z$ majore un ouvert schématiquement dense dans $X$ ; il en résulte que $Z$ est lui-même induit sur un ouvert de $X$ en vertu du lemme suivant :

\begin{lemma}[20.3.8.8]
\label{IV.20.3.8.8}
Soit $T$ un préschéma.
Alors tout sous-préschéma $W$ de $T$ qui majore un ouvert schématiquement dense $V$ de $T$ est induit sur un ouvert (schématiquement dense) de $T$.
\end{lemma}

  En effet, le sous-espace de $T$ sous-jacent à $W$ est localement fermé, donc ouvert
  